\documentclass[11pt]{article}

\usepackage[T1]{fontenc}
\usepackage[utf8]{inputenc}
\usepackage{amsmath,amssymb,amsthm}
\usepackage[margin=1in]{geometry}
\usepackage{hyperref}

\newcommand{\M}{\mathcal M}
\newcommand{\T}{\mathcal T}
\newcommand{\V}{\mathcal V}
\newcommand{\D}{\Delta}
\newcommand{\ol}[1]{\overline{#1}}

\newtheorem{theorem}{Theorem}[section]
\newtheorem{lemma}[theorem]{Lemma}
\newtheorem{proposition}[theorem]{Proposition}
\theoremstyle{remark}
\newtheorem{remark}[theorem]{Remark}

\title{Problem 9 -- GPT-6 Astra}
\author{}
\date{}

\begin{document}
\maketitle

\begin{abstract}
We prove the asserted constancy using rigidity of Teichm\"uller curves
and the boundary structure theorem for totally geodesic subvarieties.
An important preliminary observation is that the compatibility
hypothesis forces every positive-dimensional boundary locus in a
stratum with at least two nodes to remain fixed set-theoretically.
This observation also handles the possibility that a cusp of one
boundary curve lies on another positive-dimensional boundary locus.
\end{abstract}

\section{Conventions and geometric input}

We use the usual moduli-theoretic conventions: subvarieties are
algebraic, while their dependence on the disc parameter may be
holomorphic rather than algebraic. Smooth families and local
arguments on moduli spaces are understood orbifold-locally.
Thus auxiliary level structures and local labels for components and
branches of nodes may be introduced. The conclusion is a statement
on the original moduli space and is unaffected by these auxiliary
choices.

Here is one consequence of the smooth-family assumption that will
be used explicitly. Suppose that \(x\) is an isolated point of
\[
Z_0(G):=\ol{V_0}\cap\Delta_G.
\]
In a local marked chart at \(x\), the family
\[
\{(z,t):z\in Z_t(G)\},\qquad
Z_t(G):=\ol{V_t}\cap\Delta_G,
\]
has relative dimension zero and is a submersion to the parameter
disc. The holomorphic inverse function theorem therefore gives
a unique local section through the chosen marked point \(x\).
In particular, two sections through that marked point coincide
after shrinking the parameter disc. Introducing such charts is
legitimate for hypothesis~(i), which is a smooth-family hypothesis.
We will \emph{not} assume that hypothesis~(ii) holds for individual
branches in a refined or marked stratification; it will be used
only in the original moduli space.

We recall the following consequences of the boundary structure
theorem. A component of the normalization of a nodal surface will
be called \emph{active} for a boundary locus if the projection of
that locus to the corresponding pointed moduli space is
nonconstant.

\begin{theorem}[Boundary structure facts]\label{thm:structure}
Let \(N\subset\T_{g,n}\) be an algebraic totally geodesic submanifold.
The following statements hold.
\begin{enumerate}
\item
Its intersection with a stratum of augmented Teichm\"uller space
is semisimple. In particular, for a one-dimensional boundary
locus, every nonconstant factor projection is a Teichm\"uller
disc upstairs and is a local isometry on the corresponding
normalized orbifold curve downstairs.

\item
Pinchable curves have an \(N\)-equivalence relation. A pinched
multicurve is a union of entire equivalence classes. If it
contains \(r\) classes, the corresponding boundary locus has
dimension at most \(\dim N-r\). Any one class can be unpinched
while all the other classes remain pinched.

\item
For a positive-dimensional simple boundary factor, the union of
the boundary curves of its active subsurfaces is a union of
entire \(N\)-equivalence classes.
\end{enumerate}
\end{theorem}

\begin{proof}[References]
These are, respectively, Theorem~1.3; Theorem~1.6(4),
Corollary~3.17 and the dimension estimate in the proof of
Theorem~3.9; and Lemma~6.1(3) of
Arana--Herrera--Wright~\cite{AHW}. Their numbering here refers
to version~2 of that paper.
For the local-isometry assertion in the first statement, see also
Benirschke--Dozier--Rached~\cite[Theorem~1.7 and
Proposition~9.2]{BDR}.
\end{proof}

% Sources for the boundary structure and factor-projection assertions:
% Arana--Herrera--Wright and Benirschke--Dozier--Rached.
% ([arxiv.org](https://arxiv.org/html/2412.16330v2))

We need the following useful consequence of the third statement.

\begin{lemma}[Every node is seen by an active component]\label{lem:active}
Let \(W\subset\M_g\) be a totally geodesic surface, and let \(E\)
be a one-dimensional irreducible component of its intersection
with a boundary stratum. Then every node of the corresponding
stable curves is incident to at least one component that is
active for \(E\).
\end{lemma}

\begin{proof}
Choose a lift \(N\subset\T_g\) of \(W\), and a boundary locus
upstairs covering \(E\). If its pinched multicurve has \(r\)
equivalence classes, Theorem~\ref{thm:structure}(2) gives
\[
1\leq 2-r.
\]
Since this is a boundary locus, \(r\geq1\), so \(r=1\).

There is exactly one positive-dimensional simple factor, since
the locus has dimension one. Its active subsurfaces have a
nonempty union of boundary curves: the stable curve is connected,
has a node, and the locus has at least one active component.
By Theorem~\ref{thm:structure}(3), this union consists of entire
equivalence classes. There is only one such class, so the union
is the whole pinched multicurve. This is precisely the assertion.
\end{proof}

% The whole-equivalence-class assertion used in this lemma is
% Arana--Herrera--Wright, Lemma 6.1(3).
% ([public.websites.umich.edu](https://public.websites.umich.edu/~alexmw/tghhs.pdf))

\section{Deformations of boundary curves}

We will use rigidity of Teichm\"uller curves in its usual
pointed and orbifold formulation.

\begin{theorem}[Rigidity of Teichm\"uller curves]\label{thm:McM}
Let
\[
f_t:C_t\longrightarrow \M_{h,m}
\]
be a holomorphic, smoothly trivial family of orbifold maps from
finite-type curves. Suppose that \(f_0\) is a local isometry for
the complete finite-area hyperbolic metric on \(C_0\) and the
Teichm\"uller metric on the target. Then the family of pairs
\((C_t,f_t)\) is locally trivial.
\end{theorem}

\begin{proof}
For unpointed targets this is McMullen's rigidity theorem
\cite[Theorem~1]{McM}. We explain why the same argument applies
to pointed targets.

Pass to finite covers to remove orbifold stabilizers, and
identify monodromy using the smooth trivialization. For a
nonperipheral conjugacy class \(\gamma\in\pi_1(C_0)\), its
monodromy is pseudo-Anosov, with translation length equal to
the hyperbolic length \(\ell_0(\gamma)\). Holomorphic maps
contract the Teichm\"uller metric, by the
Royden--Earle--Kra theorem \cite{Royden,EK}. Consequently,
\[
\ell_t(\gamma)\geq \ell_0(\gamma)
\qquad\text{for every nonperipheral }\gamma.
\]
The marked length-spectrum domination principle used in
\cite[proof of Theorem~1, fact~(1)]{McM} implies that \(C_t\)
and \(C_0\) are conformally identical with their markings.

For completeness, the maps are then identical as marked maps.
Lift them to equivariant maps from the same hyperbolic disc to
\(\T_{h,m}\). Equality in the preceding inequality along a
hyperbolic element's axis forces that axis to map isometrically
onto the unique Teichm\"uller axis of its monodromy
\cite{Bers}. Holomorphic continuation identifies the lifted map
with the original Teichm\"uller-disc parametrization, up to an
automorphism of the source disc. Equivariance makes this
automorphism commute with the source lattice. A non-elementary
Fuchsian lattice has trivial centralizer in
\(\operatorname{PSL}_2(\mathbb R)\): commuting with two
hyperbolic elements with different axes fixes at least three
boundary points. The automorphism is therefore the identity.

Finally, the universal property of marked Teichm\"uller space
makes these identifications holomorphic in the parameter.
Thus the family of pairs is locally trivial.
\end{proof}

% Source for the rigidity theorem and the length-domination argument:
% McMullen, Rigidity of Teichmüller curves.
% ([people.math.harvard.edu](https://people.math.harvard.edu/~ctm/papers/home/text/papers/rigidity/rigidity.pdf))

\begin{lemma}[Deforming a curve in a product]\label{lem:curve}
Suppose that \(E_t\) is a holomorphic, smoothly trivial family
of algebraic curves in a product of pointed moduli spaces.
Suppose that every nonconstant factor projection of the
normalized central curve is a Teichm\"uller local isometry.
After passing to finite covers where necessary and shrinking
the parameter disc, the following hold.
\begin{enumerate}
\item
The normalized source curve can be identified with a fixed curve
\(C\), and every initially nonconstant factor map is independent
of \(t\).

\item
Every initially constant factor map remains constant on \(C\).
Its value may depend holomorphically on \(t\).

\item
Each specified point of the compactification of \(C\) determines
a holomorphic family of stable limits of \(E_t\). Its dual graph
is independent of \(t\). Every pointed component belonging to
the stable limit of an initially nonconstant factor map is fixed.
\end{enumerate}
\end{lemma}

\begin{proof}
At least one factor projection is nonconstant. Apply
Theorem~\ref{thm:McM} to that projection to identify the source
family with \(C\times\D\). Apply the same theorem to the other
initially nonconstant projections. The marked identifications
of the source agree, proving the first assertion.

Consider an initially constant projection. Its orbifold
monodromy has finite image. This image is unchanged under the
deformation. After a finite cover of \(C\) killing that image,
each deformed map lifts holomorphically to the target
Teichm\"uller space. A Bers embedding realizes that space as
a bounded domain \cite{BersEmbedding}. The coordinates of the
lift are bounded holomorphic functions on a compact Riemann
surface minus finitely many points. They extend across the
punctures and are constant. Hence the original factor map is
constant. Its value depends holomorphically on \(t\), as is seen
by evaluating at a fixed point of \(C\).

Let \(\ol C\) be the smooth compactification of \(C\).
The algebraic factor maps extend to the corresponding
Deligne--Mumford compactifications, with stable reduction or
finite covers if needed. For a nonconstant factor the extended
map is fixed; for a constant factor it is the moving smooth
point just described. Composing their product with the
clutching morphism gives the claimed family of stable limits.
Clutching and the local stable-curve description are those of
\cite{DM}. The nonconstant factors' stable limits are fixed,
and the constant factors remain smooth pointed curves of the
same topological type. Thus the dual graph does not change.
\end{proof}

\begin{lemma}[A boundary curve has a further degeneration]
\label{lem:cusp}
A positive-dimensional algebraic Teichm\"uller curve is
noncompact. Consequently, a one-dimensional boundary locus as
in Theorem~\ref{thm:structure}(1) has a stable limit in a strictly
smaller boundary stratum.
\end{lemma}

\begin{proof}
A nonzero finite-area quadratic differential generating a
Teichm\"uller disc has a cylinder; see \cite{Masur}.
Rotate the differential so that this cylinder is horizontal,
and apply Teichm\"uller flow in the direction that makes its
modulus tend to infinity. Its core curve has extremal length
tending to zero. By compactness of the thick part of moduli
space \cite{Mumford}, the resulting surfaces leave every compact
subset of moduli space. An algebraic Teichm\"uller curve is
therefore noncompact.

For a boundary curve in a product, apply this to a nonconstant
factor projection. Compactify the normalized source and the
factor maps. Some stable limit acquires an additional node,
and hence lies in a strictly smaller stratum.
\end{proof}

\section{A consequence of the compatibility hypothesis}

The next observation uses the exact, set-theoretic formulation
of hypothesis~(ii).

\begin{lemma}[Deep boundary curves cannot move]\label{lem:deep}
Let \(\Delta_G\) be a boundary stratum whose curves have at least
two nodes. For every \(t\), every positive-dimensional
irreducible component of \(Z_t(G)\) is contained in \(Z_0(G)\).
\end{lemma}

\begin{proof}
Choose a strictly larger boundary stratum \(\Delta_H\) whose
closure contains \(\Delta_G\). Such a stratum exists: smooth
one node while leaving at least one node unsmoothed.

Set
\[
U=\ol{\M}_g\setminus\ol{V_0}.
\]
For this open set, both sides of the equality in
hypothesis~(ii) are empty at \(t=0\). That hypothesis therefore
implies, for every \(t\),
\begin{equation}\label{eq:compat-empty}
Z_t(G)\cap U
=
\ol{Z_t(H)}\cap\Delta_G\cap U.
\end{equation}

Since \(\ol{V_t}\) is a surface, its boundary has dimension at
most one. Thus \(Z_t(H)\) has dimension at most one. Every
one-dimensional irreducible component of \(\ol{Z_t(H)}\) has
a dense open subset in \(\Delta_H\); its intersection with the
disjoint stratum \(\Delta_G\) is consequently finite.
Zero-dimensional components contribute no boundary points.
It follows that
\[
\ol{Z_t(H)}\cap\Delta_G
\]
is finite.

Let \(E_t\) be a positive-dimensional irreducible component of
\(Z_t(G)\). If \(E_t\) were not contained in \(\ol{V_0}\), then
\(E_t\cap U\) would be a nonempty open subset of an algebraic
curve, and hence infinite. This contradicts
\eqref{eq:compat-empty}. Therefore
\[
E_t\subset\ol{V_0}\cap\Delta_G=Z_0(G).
\]
\end{proof}

\begin{remark}\label{rem:compatibility}
The preceding proof is the only use of hypothesis~(ii).
In particular, no stronger assertion about preservation of
individual incidences between branches is being substituted
for the hypothesis in the question.
\end{remark}

\section{Proof of the assertion}

\begin{theorem}\label{thm:main}
Under the hypotheses of the question,
\[
\pi_i(\partial V_t)=\{Y_0\}
\qquad\text{for every }t\in\D.
\]
\end{theorem}

\begin{proof}
The function \(t\mapsto Y_t\) is holomorphic. Indeed, locally
choose a holomorphic section of the smooth family
\(\partial V_t\); composing with \(\pi_i\) gives \(Y_t\),
independently of the chosen section. It suffices to prove
constancy near \(0\), since the identity principle then gives
constancy on the connected disc.

Choose a one-dimensional component \(B_0\) of
\(\partial V_0\), and let \(B_t\) be its continuation in the
smoothly trivial family from hypothesis~(i).

\medskip
\noindent\textbf{Step 1: Reduction to a one-node stratum.}

If \(\Gamma\) has at least two edges,
Lemma~\ref{lem:deep} gives
\[
B_t\subset Z_0(\Gamma)=\partial V_0.
\]
Since \(\pi_i\) is equal to \(Y_t\) on \(B_t\) and equal to
\(Y_0\) on \(\partial V_0\), this already proves \(Y_t=Y_0\).

We may therefore assume that \(\Gamma\) has one edge.
It cannot be a nonseparating edge: that stratum has just one
moduli-space factor, and constancy on that factor would
contradict \(\dim B_0=1\). Thus the node is separating.
Write its vanishing curve as \(a\). Its two pointed components
will be denoted
\[
Y_t\in\M_{h,1},
\qquad
X\in\M_{g-h,1}.
\]
The \(Y\)-projection of \(B_t\) is constant, whereas the
\(X\)-projection of \(B_0\) is nonconstant.

Theorem~\ref{thm:structure}(1) and
Lemma~\ref{lem:curve} identify the normalized source of \(B_t\)
with a fixed curve and make its \(X\)-map independent of \(t\).

\medskip
\noindent\textbf{Step 2: A cusp section.}

By Lemma~\ref{lem:cusp}, choose a cusp of this fixed
\(X\)-map. Let \(X^\ast\) be its stable pointed limit.
It has at least one node. Lemma~\ref{lem:curve}(3) gives a
holomorphic section
\[
p(t)\in Z_t(G')
\]
obtained by attaching \(Y_t\) to the fixed pointed stable curve
\(X^\ast\). Here \(\Delta_{G'}\) is a fixed stratum with at least
two nodes. Put \(p_0=p(0)\).

Use local labels for the components and nodes of \(p_0\).
Along \(p(t)\), all pointed components coming from \(X^\ast\)
are fixed. Only the component \(Y_t\) can move.

The smooth locus \(Z_0(G')\) has local dimension either zero
or one at \(p_0\). We treat these possibilities separately.

\medskip
\noindent\textbf{Step 3: A positive-dimensional boundary locus
through \(p_0\).}

Suppose first that \(Z_0(G')\) has local dimension one at
\(p_0\). Hypothesis~(i) puts \(p(t)\), for small \(t\), on a
positive-dimensional component of \(Z_t(G')\).
Lemma~\ref{lem:deep} then implies
\[
p(t)\in Z_0(G').
\]
Let \(E_0\) be its local curve component through \(p_0\).

Choose any node \(b\) of \(X^\ast\). Both components incident
to \(b\), or the single component if \(b\) is a self-node,
come from \(X^\ast\); their pointed moduli are fixed along
\(p(t)\). On the other hand, Lemma~\ref{lem:active} says that
at least one of these components is active for \(E_0\).
Projection of \(E_0\) to that factor is locally injective in
a marked chart, by Theorem~\ref{thm:structure}(1).
Its value along \(p(t)\) is fixed, so \(p(t)\) itself is
constant near \(0\). In particular, \(Y_t=Y_0\).

\medskip
\noindent\textbf{Step 4: An isolated boundary point.}

It remains to suppose that \(p_0\) is isolated in \(Z_0(G')\).
Choose a marked lift of the degeneration to the closure of
a lift \(N\subset\T_g\) of \(V_0\).
Let \(A=\{a\}\), and let \(B\) be the nonempty multicurve
consisting of the additional nodes of \(X^\ast\).
Thus the nodes at \(p_0\) correspond to
\[
A\sqcup B.
\]

Because \(A\) itself occurs as the complete pinched multicurve
along \(B_0\), it is an entire \(N\)-equivalence class.
There can be at most two equivalence classes in \(A\sqcup B\),
by Theorem~\ref{thm:structure}(2) and \(\dim N=2\).
Consequently \(B\) is the other class.

Unpinch the class \(A\), leaving \(B\) pinched.
The unpinching assertion gives points arbitrarily close to
\(p_0\) in the stratum obtained by smoothing \(a\).
There is therefore a one-dimensional boundary component
\(F_0\) of \(V_0\) in that stratum with
\[
p_0\in\ol{F_0}.
\]
To justify the dimension assertion explicitly, the boundary of
the surface has dimension at most one, and a zero-dimensional
algebraic locus is finite and cannot supply points in that
stratum arbitrarily approaching \(p_0\).

% Source for the unpinching operation used in this step:
% Arana--Herrera--Wright, Lemma 3.16 and Corollary 3.17.
% ([arxiv.org](https://arxiv.org/html/2412.16330v2))

Let \(Z\) denote the component obtained by smoothing \(a\).
It contains the original \(Y_0\) and the component of
\(X^\ast\) to which \(Y_0\) was attached.
The projection of \(F_0\) to the moduli space of \(Z\) is
nonconstant: its stable limit at the specified cusp has
the node \(a\), whereas a constant smooth pointed curve
cannot have a nodal limit.

Let \(F_t\) be the continuation of \(F_0\) supplied by
hypothesis~(i). If its ambient stratum is a finite quotient
of a product, take its product cover and the corresponding
finite cover of \(F_t\). Apply Lemma~\ref{lem:curve}.
The \(Z\)-map is independent of \(t\), and its stable limit
at our specified cusp is independent of \(t\).
In particular, the pointed component inside that stable
limit which is the original \(Y_0\) remains exactly \(Y_0\).

The same lemma supplies a cusp section
\[
q(t)\in Z_t(G'),
\qquad q(0)=p_0,
\]
with the same local component and node labels.
Since \(p_0\) is isolated in the central stratum intersection,
hypothesis~(i) and the local inverse-function observation in
Section~1 imply
\[
q(t)=p(t)
\]
near \(0\), in the chosen marked chart.
Their specified \(Y\)-components agree. The \(Y\)-component
of \(q(t)\) is \(Y_0\), while that of \(p(t)\) is \(Y_t\).
Thus \(Y_t=Y_0\) near \(0\).

This completes both cases. Descending any auxiliary markings
or finite covers and applying the identity principle yields
\[
\pi_i(\partial V_t)=\{Y_0\}
\qquad(t\in\D).
\]
\end{proof}

\begin{remark}
Only the central surface and its central boundary curves were
assumed to be totally geodesic. For \(t\ne0\), rigidity of the
nonconstant projections of boundary curves follows from
Theorem~\ref{thm:McM}; total geodesicity of \(V_t\) is never
used.
\end{remark}

\begin{thebibliography}{99}

\bibitem{AHW}
F.~Arana--Herrera and A.~Wright,
\emph{The geometry of totally geodesic subvarieties of moduli
spaces of Riemann surfaces},
arXiv:2412.16330v2, 2025.

\bibitem{BDR}
F.~Benirschke, B.~Dozier, and J.~Rached,
\emph{The boundary of a totally geodesic subvariety of moduli
space},
arXiv:2412.06765.

\bibitem{Bers}
L.~Bers,
\emph{An extremal problem for quasiconformal mappings and a
theorem by Thurston},
Acta Math. \textbf{141} (1978), 73--98.

\bibitem{BersEmbedding}
L.~Bers,
\emph{Spaces of Riemann surfaces as bounded domains},
Bull. Amer. Math. Soc. \textbf{66} (1960), 98--103.

\bibitem{DM}
P.~Deligne and D.~Mumford,
\emph{The irreducibility of the space of curves of given genus},
Publ. Math. Inst. Hautes \'Etudes Sci. \textbf{36} (1969),
75--109.

\bibitem{EK}
C.~J.~Earle and I.~Kra,
\emph{On isometries between Teichm\"uller spaces},
Duke Math. J. \textbf{41} (1974), 583--591.

\bibitem{Masur}
H.~Masur,
\emph{Closed trajectories for quadratic differentials with an
application to billiards},
Duke Math. J. \textbf{53} (1986), 307--314.

\bibitem{McM}
C.~T.~McMullen,
\emph{Rigidity of Teichm\"uller curves},
Math. Res. Lett. \textbf{16} (2009), no.~4, 647--650.

\bibitem{Mumford}
D.~Mumford,
\emph{A remark on Mahler's compactness theorem},
Proc. Amer. Math. Soc. \textbf{28} (1971), 289--294.

\bibitem{Royden}
H.~L.~Royden,
\emph{Automorphisms and isometries of Teichm\"uller space},
in \emph{Advances in the Theory of Riemann Surfaces},
Annals of Mathematics Studies, vol.~66,
Princeton University Press, 1971, 369--383.

\end{thebibliography}

\end{document}